%% kapa: exam=1 round=25 year=2014 period=1 form=A subject=civil_law subject_ko=민법 seq=30 number=30
%% source: Q-Net 감정평가사 기출문제 (cst00309) artlSeq=5101207 file=extracted/1008/5101207/2014년 제25회 감정평가사 1차시험 기출문제/1교시(A형).hwp
%% conversion: kordoc 4.24.51 -> tools/gen1.py (auto); review=proofread-sonnet; flags=none
\begin{question}{30}
집합건물의 구분소유에 관한 설명으로 옳지 않은 것은? (다툼이 있으면 판례에 의함)
\begin{choices}
  \item 구분소유자는 전유부분의 보존․개량을 위하여 다른 구분소유자의 전유부분의 사용을 청구할 수 있다.
  \item 구조상․이용상의 독립성을 갖춘 건물부분은 구분소유권의 목적이 될 수 있다.
  \item 구분소유자의 대지사용권은 규약으로 달리 정함이 없는 한, 그가 가지는 전유부분의 처분에 따른다.
  \item 아파트의 지하실은 특별한 사정이 없는 한, 구분소유권의 목적이 될 수 없다.
  \item 각 구분소유자의 공용부분에 대한 지분은 규약으로 달리 정함이 없는 한, 그가 가지는 전유부분의 가액 비율에 따른다.
\end{choices}
\end{question}
\answer{5}
